% Post-checkpoint-14 extension of the finite-puncture reduction.
% Complete author/root/B/C actual-file and scoped primary audits passed.
\begin{theorem}[Proved: surface reduction with controlled exceptional centres]
\label{D:generic-surface-controlled-boundary}
Let $k$ be algebraically closed, let $S$ be a normal proper integral
$k$-scheme of dimension $d\ge2$, and let $T\subsetneq S$ be closed.
Write $U=S\setminus T$. Suppose that there is a proper birational
morphism
\[
                      q:Y\longrightarrow S
\]
with $Y$ normal and projective over $k$, satisfying this condition:
every prime divisor $E\subset Y$ whose image has codimension at least
two in $S$ and is contained in $T$ has zero-dimensional image.

Then there are an algebraically closed extension $K/k$ of finite
transcendence degree, a normal proper integral surface $X/K$, and a
finite morphism $g:X\to S_K$ onto its image, such that, with
$T_X=g^{-1}(T_K)$, the induced map on the original-boundary tame
quotients is surjective:
\[
                 Q(X,T_X)\twoheadrightarrow Q(S,T).
                 \tag{D.CB.1}
\]
Here tameness is required only at the original codimension-one boundary
points in each of the indicated compactifications.

If $T$ has codimension at least two in $S$, then $T_X$ is finite and
there is a surjection of full fundamental groups
\[
                 \pi_1(X\setminus T_X)\twoheadrightarrow\pi_1(U).
                 \tag{D.CB.2}
\]
Thus, in characteristic $p>0$, all the already proved surface bounds
for all-field $H^1$, defining-characteristic quadratic representation
counts, and bounded-scalar projective representation counts hold for
$Q(S,T)$ under the displayed geometric condition. Its maximal
prosolvable quotient and the previously specified bounded-scalar
simple-product quotients are finitely generated. For (D.CB.2) these
conclusions concern the full group $\pi_1(U)$. We assert neither an
$H^2$ transfer nor full finite generation of these groups.
\end{theorem}

\begin{proof}
When $d=2$ the assertion uses $X=S$, $K=k$, so suppose $d\ge3$.
Let $D_1,\ldots,D_s$ be the codimension-one components of $T$ and
let $B_j\subset Y$ be their strict transforms. A proper birational
morphism to a normal scheme is an isomorphism over an open containing
all codimension-one points. Thus every divisorial component of
$q^{-1}(T)$ is either one of the $B_j$ or an exceptional divisor whose
image has codimension at least two. By hypothesis the latter divisors
all have point images. Denote their union by $E_0$ and the finite set
of their images by $P$. The remaining components of $q^{-1}(T)$ have
dimension at most $d-2$
\citep[Tag 0BFP, Lemma 33.17.3 and proof]{StacksDControlledBoundary}.

Choose a dense smooth open $W\subset U$ over which $q$ is an
isomorphism and an embedding $Y\hookrightarrow\mathbf P^N_k$.
Perform the same fixed tower of $d-2$ geometric generic hyperplane
sections used in Theorem~\ref{D:generic-surface-finite-punctures}.
Its final intersection $W_*$ is a smooth integral surface. Every
original connected finite etale cover of $U$ restricts to a connected
cover of $W_*$, by the exact finite-cover Bertini argument in that
theorem.

We also impose the following finite list of geometric generic
conditions. A closed subscheme of $Y$ of dimension at most $d-2$
meets these $d-2$ generic hyperplanes in dimension at most zero.
Indeed at each stage a generic hyperplane contains no
positive-dimensional irreducible component of the preceding cut;
such components drop dimension by one, and a generic hyperplane
avoids all zero-dimensional components. Apply this to the
nondivisorial components of $q^{-1}(T)$.

There is an open $\Omega\subset S$ containing the generic points of
all the $D_j$ such that $q$ is an isomorphism over $\Omega$, the scheme
$\Omega$ is smooth, and the boundary there is a disjoint union of
smooth Cartier divisors $D_j\cap\Omega$; all other components of $T$
are absent. To obtain it, remove the non-isomorphism locus of $q$,
the singular locus of $S$, the singular and pairwise intersection
loci of the $D_j$, and the remaining boundary components, shrinking
only by subsets which avoid each generic point of $D_j$.
For each $j$, the closed subset
\[
                         B_j\setminus q^{-1}(\Omega)
\]
has dimension at most $d-2$. Its intersection with the generic surface
is therefore finite. Apply smooth Bertini successively also to
$q^{-1}(\Omega)$ and to its smooth divisors $B_j\cap q^{-1}(\Omega)$.
The surface section is smooth there and meets each such divisor
transversely at the generic points of its boundary curves. These are
finitely many additional conditions on the same fixed generic tower,
and follow from Stacks, Lemma 33.47.3
\citep[Tag 0FD6]{StacksDGenericSurface}.

Let $Z$ be the reduced closure of $W_*$ in $Y_K$, let $C=q_K(Z)$
with its reduced induced structure, and let $X=C^\nu$. Exactly as in
Theorem~\ref{D:generic-surface-finite-punctures}, $C$ is an integral
proper surface, its normalization is finite, and $X$ contains $W_*$
as a dense open. On $q_K^{-1}(\Omega_K)$ the complete-intersection
section is smooth of pure dimension two; its closed pure
two-dimensional subscheme $Z\cap q_K^{-1}(\Omega_K)$ is a union of
irreducible components, and is therefore smooth too. Consequently
\[
                  \pi_1(X\setminus T_X)\twoheadrightarrow\pi_1(U)
                  \tag{D.CB.3}
\]
is surjective: every original connected cover remains connected on
$W_*$, and each component of its finite etale pullback to the normal
integral scheme $X\setminus T_X$ meets that open.

We describe the codimension-one boundary of $X$. The image of
$Z\cap E_{0,K}$ lies in the finite set $P_K$. Every intersection of
$Z$ with a remaining nondivisorial component of $q^{-1}(T)_K$ is
finite. Hence each curve in $C\cap T_K$ comes from a curve in
$Z\cap B_{j,K}$ for some $j$. Its generic point lies over $\Omega_K$,
where $q_K$ is an isomorphism and $Z$ is smooth and transverse to
$B_{j,K}$. Normalization of $C$ changes nothing near that point.
Thus every codimension-one boundary point of $X$ is of precisely this
transverse original-divisor type. If there are no $D_j$, the entire
set $C\cap T_K$ is finite and so is $T_X$; this proves the extra
boundary assertion in (D.CB.2).

Now let $V\to U$ be a connected original-boundary tame cover. Over
$\Omega$ it is a finite etale cover of the complement of the displayed
smooth divisors, tame at their generic points. The local absolute
Abhyankar lemma gives its precise description: near each such
boundary point, after adjoining an $n$th root of a local boundary
parameter $t$, with $n$ prime to the characteristic exponent, its
normalization is finite etale. Equivalently over a strict localization
the cover is dominated by that Kummer cover
\citep[Expose XIII, Propositions 5.1--5.2 and Corollary 5.3,
pp.~428--430]{C_SGA1}.
After any algebraically closed field extension this description
persists. At a generic boundary point of $X$ the restriction of $t$
is a uniformizer because of transversality. Pulling back the Kummer
cover and its finite etale extension therefore shows that the
pullback of $V$ is tame at this DVR, including separability of the
residual extension. This applies to every codimension-one boundary
point of $X$. No condition is imposed at its isolated boundary
points.

Thus every finite quotient of $Q(S,T)$ pulls back to a tame cover
on $(X,T_X)$. The composite of (D.CB.3) with the quotient to $Q(S,T)$
annihilates the wild inertia groups defining $Q(X,T_X)$, and hence
factors through (D.CB.1). It remains surjective. If the original
boundary has codimension at least two, there are no original
divisorial tameness conditions and (D.CB.3) is exactly (D.CB.2).

The quantitative consequences follow by degree-one inflation and
injectivity of representation classes under surjective precomposition,
using Theorem~\ref{D:tame-surface-all-characteristics-H1},
Theorem~\ref{D:tame-global-quadratic-representation-bound}, and the
surface bounded-scalar results. The argument is the direct transfer
in Corollary~\ref{D:normal-finite-puncture-quantitative-transfer}; it
does not use degree-two inflation.
\end{proof}

\begin{remark}[Exact geometric limitation]
The hypothesis concerns the chosen projective modification and only
its exceptional divisors whose centres are contained in the omitted
set. It holds, for example, if the modification has no exceptional
divisors over that set except ones contracted to finitely many points.
It does not say that an arbitrary proper scheme admits such a
modification for every positive-dimensional boundary. In the blowup
of $\mathbf P^3$ in a line from
Example~\ref{D:Chow-surface-boundary-obstruction}, the exceptional
divisor maps onto the omitted line, so precisely this condition fails.
We have not converted the general positive-dimensional-boundary case
to finite punctures on a surface.
\end{remark}
